\chapter{الأعداد حتى 100}\label{ch:g1:tens}

عدّ كومة كبيرة واحدًا واحدًا يطول بلا نهاية. والحيلة الكبرى: اِربط
الأشياء في حزم من \emph{عشرة}، ثم عُدّ الحزم. هكذا تُبنى الأعداد
حتى $100$ --- ولهذا تُكتب
برقمين.

\section{حزم العشرة}

\begin{definition}[العشرات والآحاد]\label{def:g1:tens:def}
\emph{العشرة} حزمة فيها عشر وحدات. والعدد $34$ يعني $3$
عشرات و $4$ آحاد:
\[
34 = 10 + 10 + 10 + 4 .
\]
الرقم الأيسر يعدّ الحزم، والرقم الأيمن يعدّ الآحاد المفردة.
\end{definition}

\begin{omfigure}
\begin{tikzpicture}[scale=0.75]
  % three bundles of ten sticks, each tied with a rounded band
  \foreach \b in {0,1,2} {
    \begin{scope}[xshift=\b*2.4cm]
    \foreach \i in {0,...,9} \draw[omDef, very thick]
      ({0.16*\i},0) -- ({0.16*\i},1.7);
    \draw[omExo, thick, rounded corners=3pt]
      (-0.22,0.62) rectangle (1.66,1.05);
    \end{scope}
  }
  % four loose sticks
  \foreach \i in {0,...,3} \draw[omThm, very thick]
    (7.7+0.4*\i,0) -- (7.7+0.4*\i,1.7);
  \node[right, font=\small] at (9.7,0.85)
    {$3$ حزم و $4$ أعواد: $34$};
\end{tikzpicture}

{\small أربعة وثلاثون عودًا: ثلاث حزم من عشرة (كل حزمة مربوطة
بشريط) وأربعة أعواد مفردة. لا حاجة إلى العدّ حتى $34$ واحدًا واحدًا!}
\end{omfigure}

\begin{example}[قراءة الأعداد]\label{ex:g1:tens:reading}
$20$ عشرتان (عشرون)، و $30$ ثلاث عشرات (ثلاثون)\,\dots\ حتى
$90$ (تسعون) و $100$ (مئة: عشر عشرات!). وبينها:
$47$ يُقرأ «سبعة وأربعون» --- أربعون لأجل $4$ عشرات، وسبعة
للآحاد.
\end{example}

\begin{omfigure}
\begin{tikzpicture}[scale=0.9]
  \draw (0,0) grid[xstep=2.2, ystep=0.8] (4.4,1.6);
  \node[font=\small] at (1.1,1.2) {عشرات};
  \node[font=\small] at (3.3,1.2) {آحاد};
  \node[font=\large, omDef] at (1.1,0.4) {$3$};
  \node[font=\large, omThm] at (3.3,0.4) {$4$};
\end{tikzpicture}

{\small جدول المراتب للعدد $34$: الرقم $3$ في مرتبة العشرات،
والرقم $4$ في مرتبة الآحاد. وهذا الجدول الصغير يزداد عمودًا كل
سنة!}
\end{omfigure}

\begin{example}[الأرقام نفسها، وأعداد مختلفة]\label{ex:g1:tens:place}
$27$ و $72$ يستعملان الرقمين نفسيهما لكنهما مختلفان جدًّا: فالعدد
$27$ فيه $2$ عشرات و $7$ آحاد، والعدد $72$ فيه $7$ عشرات و $2$ آحاد
--- أكبر بكثير! \emph{مرتبة} الرقم هي التي تقرر ما يعدّه.
\end{example}

\section{شبكة المئة}

\begin{omfigure}
\begin{tikzpicture}[scale=0.52]
  \foreach \r in {0,...,9} {
    \foreach \c in {0,...,9} {
      \pgfmathtruncatemacro{\n}{\r*10+\c+1}
      \draw[gray!50] (\c,-\r) +(-0.5,-0.5) rectangle +(0.5,0.5);
      \node[font=\footnotesize] at (\c,-\r) {\n};
    }
  }
  \draw[omThm, very thick] (3,-4) +(-0.5,-0.5) rectangle +(0.5,0.5);
  \draw[omDef, very thick] (3,-5) +(-0.5,-0.5) rectangle +(0.5,0.5);
\end{tikzpicture}

{\small شبكة المئة: كل سطر عشرة. وخطوة واحدة نحو اليمين تعني
$+1$؛ وخطوة واحدة نحو \emph{الأسفل} تعني $+10$ (من $44$ الأحمر
نزولًا إلى $54$ الأزرق).}
\end{omfigure}

\begin{method}[التنقّل على الشبكة]\label{met:g1:tens:grid}
\begin{enumerate}
  \item $+1$ / $-1$: خطوة واحدة يمينًا / يسارًا؛
  \item $+10$ / $-10$: خطوة واحدة نزولًا / صعودًا --- ولا يتغير
        إلا رقم العشرات: $37 + 10 = 47$؛
  \item العدّ بالعشرات: $10, 20, 30, \dots$ --- اتبع العمود
        الأخير من الشبكة نزولًا.
\end{enumerate}
\end{method}

\begin{example}[العدّ اثنين اثنين وخمسة خمسة]\label{ex:g1:tens:skip}
القفز بمقدار $2$: $2, 4, 6, 8, 10, 12, \dots$ (الأعداد الزوجية).
والقفز بمقدار $5$: $5, 10, 15, 20, 25, \dots$ --- وهي تنتهي بالرقم $5$ أو
$0$ وترسم شكلًا جميلًا على الشبكة. والعدّ بالقفز يمهّد للضرب
(\cref{ch:g2:mult}).
\end{example}

\section{مقارنة الأعداد الأكبر}

\begin{method}[المقارنة بالعشرات]\label{met:g1:tens:compare}
صاحب العشرات الأكثر يفوز: $61 > 58$ لأن $6$ عشرات تغلب $5$
عشرات --- مع أن $8$ أكثر من $1$! ولا تقرر الآحاد إلا إذا تساوت
العشرات: $43 < 47$.
\end{method}

\begin{example}\label{ex:g1:tens:compare}
رتّب $35$ و $53$ و $29$: بالعشرات أولًا ($2$ و $3$ و $5$):
\[
29 < 35 < 53 .
\]
\end{example}

\section{تمارين}

\begin{exercise}[$\star$]\label{exo:g1:tens:1}
كم عشرة وكم واحدًا في: $26$؛ \ $40$؛ \ $73$؛ \ $99$؟
\end{exercise}

\begin{exercise}[$\star$]\label{exo:g1:tens:2}
اُكتب العدد: $5$ عشرات و $2$ آحاد؛ \ $8$ عشرات؛ \ $1$ عشرة
و $9$ آحاد؛ \ $10$ عشرات.
\end{exercise}

\begin{exercise}[$\star$]\label{exo:g1:tens:3}
اُكتب بالأرقام: ثلاثة وستون؛ \ أربعون؛ \ خمسة وثمانون؛ \
مئة.
\end{exercise}

\begin{exercise}[$\star$]\label{exo:g1:tens:4}
احسب بالتنقّل على شبكة المئة: $23 + 10$؛ \quad $67 + 10$؛
\quad $45 - 10$؛ \quad $80 - 10$.
\end{exercise}

\begin{exercise}[$\star$]\label{exo:g1:tens:5}
أكمل: $10, 20, 30, ?, ?, ?$. \quad
$5, 10, 15, ?, ?, ?$. \quad $12, 14, 16, ?, ?, ?$.
\end{exercise}

\begin{exercise}[$\star$]\label{exo:g1:tens:6}
اُنقل وأكمل بالإشارة $<$ أو $>$:
\[
34 \;?\; 43, \qquad
58 \;?\; 61, \qquad
70 \;?\; 69, \qquad
99 \;?\; 100 .
\]
\end{exercise}

\begin{exercise}[$\star$]\label{exo:g1:tens:7}
رتّب من الأصغر إلى الأكبر: $46$؛ \ $19$؛ \ $64$؛ \ $91$؛ \
$41$.
\end{exercise}

\begin{exercise}[$\star$]\label{exo:g1:tens:8}
يربط بستانيّ $52$ زهرة في باقات من عشرة. كم باقة كاملة تخرج، وكم
زهرة تبقى مفردة؟
\end{exercise}

\begin{exercise}[$\star\star$]\label{exo:g1:tens:9}
على شبكة المئة، اِبدأ من $36$؛ ثم اِنزل خطوتين واذهب خطوة واحدة
يسارًا. أين تصل؟ وأي عملية حسبتها للتوّ؟
\end{exercise}
